% topic: algorithms
\question พิจารณาขั้นตอนวิธีซึ่งแสดงด้วยรหัสเทียม (pseudocode) ต่อไปนี้
\begin{pseudocode}
\begin{enumerate}
    \item รับข้อมูลเป็นสตริง $S$ ความยาว 4 ตัวอักษร ที่ประกอบด้วยตัวอักษร A, B, C, D
    \item รับข้อมูลเป็นจำนวนเต็มบวกเก็บไว้ในตัวแปรชื่อ $n$
    \item ทำซ้ำจำนวน $n$ รอบ โดยในแต่ละรอบ ให้ปรับเปลี่ยนสตริง $S$ ตามลำดับขั้นตอนต่อไปนี้
    \begin{enumerate}
        \item สลับตำแหน่งตัวอักษร โดยนำตัวอักษรในตำแหน่งที่ 2 ไปไว้หน้าสุด แล้วตามด้วยตัวอักษรในตำแหน่งที่ 4, 1 และ 3 ตามลำดับ (เช่น หากสตริงคือ 1-2-3-4 ผลลัพธ์หลังสลับจะเป็น 2-4-1-3)
    \end{enumerate}
\end{enumerate}
\end{pseudocode}
ถ้าป้อนข้อมูลเริ่มต้นเป็นสตริง ABCD และจำนวนเต็ม $n = 2026$ หลังการทำงานครบ $n$ รอบ สตริงผลลัพธ์จะมีค่าตรงกับข้อใด \\
\choiceshorizontal{BDAC}{DCBA}{CADB}{ABCD}
% answer: 2